\documentclass[10pt,a4paper]{amsart}
\usepackage[margin=19mm]{geometry}
\usepackage{amsmath,amssymb,mathtools}
\usepackage[scr=rsfs,cal=euler]{mathalfa}
\usepackage{xfrac}
\usepackage[unicode,hidelinks,bookmarks=false]{hyperref}
\newcommand{\ie}{i.e.\ }
\newcommand{\cf}{cf.\ }
\newcommand{\resp}{respectively\ }
\newcommand{\Subsection}{\subsection}
\providecommand{\nobreakdash}{\nobreak\hbox{-}}
\setlength{\emergencystretch}{2em}
\multlinegap0pt
\usepackage{float}
\usepackage{amssymb}
\usepackage{tikz}
\usepackage{xcolor}
\usepackage[normalem]{ulem}
\newtheorem{theorem}{Theorem}
\newtheorem{lemma}{Lemma}
\newcommand{\N}{{\mathbb{N}}}
\newcommand{\R}{{\mathbb{R}}}
\newcommand{\Z}{{\mathbb{Z}}}
\title{Persistence of \(n\)-Species Lotka--Volterra Models with Periodic Pulses}
\author{Eleanor Courcelle, Jane Shaw MacDonald, Swati Patel}
\date{Source: arXiv:2609.09343v1 (CC BY 4.0)}
\begin{document}
\maketitle

\noindent\emph{Source and licence.} Persistence of \(n\)-Species Lotka--Volterra Models with Periodic Pulses. Version arXiv:2609.09343v1 (CC BY 4.0), distributed by arXiv under the Creative Commons Attribution 4.0 International licence, http://creativecommons.org/licenses/by/4.0/, as recorded in the arXiv licence metadata for this exact version and shown on the abstract page https://arxiv.org/abs/2609.09343v1. Full licence notice: \texttt{licenses/arxiv-2609.09343-CC-BY-4.0.txt}.\par\smallskip

\noindent\textit{Editorial scope.} Counted unit: the article's theorem eqxn:lv\_cancer\_model\_permanence\_req\_1 (source0.tex lines 638--651). The article's complete original proof of it is delivered with it (source0.tex lines 652--704, one \texttt{proof} environment of 53 lines). No mathematics is written, repaired or re-derived: the statement and the proof are the article's own lines, in the article's own notation, and no retained line is substituted or re-flowed.  Retained uncounted as the source-local context the proof needs: lines 100--105; lines 141--143; lines 214--220; lines 216--218; lines 538--540; lines 555--561; lines 566--571; lines 574--584; lines 588--597; lines 590--592; lines 593--595.  Nothing was omitted from any retained span, so the package carries no omission record.  No retained line contains a citation, so the wrapper carries no bibliography and no bibliography entry.  No retained line contains a figure environment, an image-inclusion command or drawing code, so no image is delivered and the package declares no asset dependency.  Editorial formatting only, in addition: an AMS \texttt{amsart} preamble that restates the article's own declarations and macros used by the retained lines, namely the environments \texttt{theorem}, \texttt{lemma} on separate counters, together with the standard packages those lines need.  The retained environments print in this document as Theorem 1, Lemma 1 in order of appearance, whereas the article prints the same items with the numbering of its own full document.  The complete native source of the article as distributed in the arXiv e-print for this exact version is bundled unchanged as \texttt{sources/209803/source0.tex}.
\begin{equation} \label{eqxn:pulsed_model}
    \begin{split}
        \frac{dx_i}{dt} = x_i(t) f_i(x) ,\;  i = 1 \ldots n\; ,\ t \neq k\tau; k \in \Z_{+}\\
        x_i(k\tau^+) = (1+h_i)x_i(k\tau^-)
    \end{split}
\end{equation}

\begin{equation}\label{discretemap}
x_{k+1}= r(\Phi(x_k, \tau))   
\end{equation}

\begin{theorem} \label{thm:main}
  Let $\mathcal{M}_\pi = \{M_1, M_2,... M_l\}$ be a Morse Decomposition for $E \cap \Gamma$ with respect to (\ref{discretemap}). If, for each $M_k\in \mathcal{M}$ there exists $p_{k1},..., p_{kn} > 0$ such that for every $x \in M_k$ there is an $\ell_x \in \mathbb{N}$ satisfying
    \begin{equation}\label{eqxn:persistence_inequality_pulsed}
        \sum^n_{i=1} p_{ki} \sum^{\ell_x}_{j=0} \int_0^\tau f_i(\Phi(\pi(x,j),s)) + \frac{\ln{(1+h_i)}}{\tau} ds > 0
    \end{equation}
    then (\ref{eqxn:pulsed_model}) is permanent.
\end{theorem}

    \begin{equation}
        \sum^n_{i=1} p_{ki} \sum^{\ell_x}_{j=0} \int_0^\tau f_i(\Phi(\pi(x,j),s)) + \frac{\ln{(1+h_i)}}{\tau} ds > 0
    \end{equation}

\begin{theorem} \label{thm:robust_permanence}
    If (\ref{eqxn:pulsed_model}) satisfies condition (\ref{eqxn:persistence_inequality_pulsed}) for permanence, then (\ref{eqxn:pulsed_model}) is also robustly permanent. 
\end{theorem}

\begin{equation}\label{eqxn:lv_cancer_model}
\begin{split}    
\frac{d x_s}{d t} = x_s(a_s - x_sb_{ss} - x_rb_{sr}), \; t \neq k\tau, k \in \N \\
\frac{dx_r}{dt} = x_r(a_r - x_sb_{rs} - x_rb_{rr}),  \; t = k\tau\\
x_i(k\tau)= (1 + h_i)x_i(k\tau^-). 
\end{split}
\end{equation}

    \begin{equation}\label{eqxn:lv_1d}
    \begin{split}
        \frac{dx}{dt} = x(a - b x), \; t \neq k\tau; k \in \Z_+ \\
        x(k\tau) = (1+h)x(\tau^-)
    \end{split}
    \end{equation}

\begin{lemma}[Jin et al.; Lemma 2.4]\label{lemma:jin_1d_periodic_soln}
    Consider (\ref{eqxn:lv_1d}). If
    \begin{equation*}
        \ln{(1+h) + \tau a > 0}
    \end{equation*}
    then (\ref{eqxn:lv_1d}) has a unique $\tau$-periodic solution $x^*(t)$, and $x^*(t)$ is globally, asymptotically stable in the sense that $\lim_{t \to \infty}|x(t) - x^*(t)| = 0$, where $x(t)$ is any solution of system (\ref{eqxn:lv_1d}) with positive initial value $x(0) > 0$. Conversely, if
    \begin{equation*}
        \ln{(1+h) + \tau a < 0}
    \end{equation*}
    then $\lim_{t \to \infty}|x(t)| = 0$ where $x(t)$ is any solution of system (\ref{eqxn:lv_1d}) with positive initial value $x(0) > 0$.
\end{lemma}

\begin{lemma}\label{lemma:periodic_integral_of_lv_1d}
    Consider (\ref{eqxn:lv_1d}). If there exists a nontrivial $\tau$-periodic solution $x^*$ for (\ref{eqxn:lv_1d}) such that $x^*(t + \tau) = x^*(t)$, then the following statements are true for all solutions $x(t)$ with initial conditions in $x^*$.
    \begin{equation}\label{eqxn:periodic_integral_of_f}
        \int_0^\tau f(x) \; dt = -\ln{(1 + h)}
    \end{equation}
    \begin{equation}\label{eqxn:periodic_integral_of_x}
        \int_0^\tau x(t) \; dt = \frac{a\tau  + \ln{(1+h)}}{b} 
    \end{equation}
           
\end{lemma}

    \begin{equation}
        \int_0^\tau f(x) \; dt = -\ln{(1 + h)}
    \end{equation}

    \begin{equation}
        \int_0^\tau x(t) \; dt = \frac{a\tau  + \ln{(1+h)}}{b} 
    \end{equation}

\begin{theorem}
    Consider (\ref{eqxn:lv_cancer_model}). Assume
    \begin{equation}\label{eqxn:lv_cancer_model_permanence_req_1}
        \tau a_s + \ln{(1 + h_s)} > 0 , \; \tau a_r + \ln{(1 + h_r)} > 0
    \end{equation}
    If 
    \begin{equation}\label{eqxn:lv_cancer_model_permanence_req_2}
        \begin{split}
            a_s\tau + \ln{|1+h_s|} > \left(a_r\tau + \ln{|1 + h_r|}\right)\frac{b_{sr}}{b_{rr}}\\
            a_r\tau + \ln{|1+h_r|} > \left(a_s\tau + \ln{|1 + h_s|}\right)\frac{b_{rs}}{b_{ss}}
        \end{split}
    \end{equation}
    then the system is robustly permanent.
\end{theorem}

\begin{proof}
    Let $\pi$ be the discrete time map for (\ref{eqxn:lv_cancer_model}) defined as in (\ref{discretemap}).
    
    We first construct an appropriate Morse decomposition $\mathcal{M}_\pi$. The extinction set $E$ of (\ref{eqxn:lv_cancer_model}) is the set $\{(x_r,x_s)| x_rx_s=0\}$. By (\ref{eqxn:lv_cancer_model_permanence_req_1}) and Lemma \ref{lemma:jin_1d_periodic_soln}, there are unique $\tau$-periodic solutions $x_r^*(t)$ and $x_s^*(t)$ in the sets $\{(0, x_r)|x_r>0\}$ and $\{((x_s,0)|x_s>0\}$, respectively.  From these periodic solutions, we define the sets $M_1$ and $M_2$ in $E$. Then, if $\Gamma$ is the global attractor on (\ref{eqxn:lv_cancer_model}), a valid Morse Decomposition for $E\cap\Gamma$ is  
    $$
    \mathcal{M}_\pi = \{ M_1, M_2, M_3 \}
    $$ with
    %$$
    %M_1 = \{(0,x_s^*(\tau))\}, M_2 =\%{(x_s^*(\tau),0)\}, M_3 =  \{(0,0)\}
   % $$
    where $M_3=\{(0.0)\}$.
    Now, for each $k \in \{1, 2, 3\}$ and $M_k \in \mathcal{M}$ we consider the existence of $p_{k1},p_{k2} >0$ satisfying condition (\ref{eqxn:persistence_inequality_pulsed}).

    \textit{($k = 1$)} 
    Let  $z \in M_1$, so $z = (0,x_s^*(\tau))$. Recall $x_s^*(t)$ is a $\tau$-periodic solution of (\ref{eqxn:lv_cancer_model}) so for all $s \in \R$,
    $$
    \Phi(\pi(z,0),s) = \Phi(\pi(z,j),s), \; j\in \N
    $$ By Lemma \ref{lemma:periodic_integral_of_lv_1d}.\ref{eqxn:periodic_integral_of_f} we have
    $$
    \int_0^\tau f_s(\Phi (\pi(z,j),s)) + \frac{\ln{(1+h_s)}}{\tau} \; dt = -\ln{(1+h_s)} + \ln{(1+h_s)} = 0
    $$
    Likewise, expanding $f_r$ and applying Lemma \ref{lemma:periodic_integral_of_lv_1d}.\ref{eqxn:periodic_integral_of_x} we obtain
    $$
    \int_0^\tau a_r + b_{rs}x_s(t) +b_{rr}x_r(t) \; dt = a_r\tau + b_{rs}\left(-\frac{a_s\tau + \ln{(1+h_s)}}{b_{ss}}\right) + \ln{(1 + h_r)}
    $$ So combined we see that the following expressions are equivalent
    \begin{equation*}
        \begin{split}
            p_{1s} \int_0^\tau f_s(\Phi (\pi(z,j),s)) + \frac{\ln{(1+h_s)}}{\tau} \; ds + p_{1r} \int_0^\tau f_r(\Phi (\pi(z,j),s)) + \frac{\ln{(1+h_r)}}{\tau} \; ds=\\
            p_{1r} \left[a_r\tau + b_{rs}\left(-\frac{a_s\tau + \ln{(1+h_s)}}{b_{ss}}\right) + \ln{(1 + h_r)} \right]\\
        \end{split}
    \end{equation*}
    Therefore, since $a_r\tau + \ln{(1 + h_r)} > \left(a_s\tau + \ln{|1 + h_s|}\right)\frac{b_{rs}}{b_{ss}}$, for $p_{1s}, p_{1r} = 1$, inequality (\ref{eqxn:persistence_inequality_pulsed}) is satisfied for all $z \in M_1$. \\
    
    \textit{($k = 2$)}  
    By symmetry, we may interchange $x_s$ and $x_r$ in the $k = 1$ case above to see that $a_s\tau + \ln{|1+h_s|} > \left(a_r\tau + \ln{|1 + h_r|}\right)\frac{b_{sr}}{b_{rr}}$ implies
    \begin{equation*}
        p_{2s}\int_0^{\tau} f_s(\Phi(\pi(z,j),s)) + \frac{\ln{(1+h_s)}}{\tau}\; ds +  p_{2r} \int_0^{\tau} f_r(\Phi(\pi(z,j),s)) + \frac{\ln{(1+h_r)}}{\tau}\; ds > 0
    \end{equation*}
    is true for $p_{2s}, p_{2r} = 1$ and any $z \in M_2$.

    \textit{($k = 3$)}
    Expanding $f_s(z), f_r(z)$ with $z = (0,0)$ gives $f_i(z) = a_i$. Hence, we have
    \begin{equation*}
    \begin{split}
        p_{3s} \int_0^{\tau} f_s(\Phi(\pi(z,j), s)) + \frac{\ln{(1+h_s)}}{\tau}\; ds +  p_{3r} \int_0^{\tau} f_r(\Phi(\pi(z,j), s)) + \frac{\ln{(1+h_r)}}{\tau}\; ds \\
        =
        p_{3s} (a_s \tau + \ln{(1+h_s)}) + p_{3r}( a_r \tau + \ln{(1+h_r)})
    \end{split}
    \end{equation*}
    So by (\ref{eqxn:lv_cancer_model_permanence_req_1}) for $p_{3s}, p_{3r} = 1$, (\ref{eqxn:persistence_inequality_pulsed}) holds for $z = (0,0)$. 
    
    Therefore for each $M_k \in \mathcal{M}$, there are $p_{ks}, p_{kr} > 0$ satisfying (\ref{eqxn:persistence_inequality_pulsed}), and by Theorem \ref{thm:main} and Theorem \ref{thm:robust_permanence}, (\ref{eqxn:lv_cancer_model}) is robustly permanent.
\end{proof}

\end{document}
